% Source: Weinberg GR, physical PDF pp. 433--436
%         (printed pp. 412--415).
% Coverage: Section 14.2, The Robertson--Walker Metric; equations
%           (14.2.1)--(14.2.21).
% Figures/tables/footnotes: one source footnote following (14.2.21); no
%                          figures, tables, or reference markers.
% Status: source-reviewed and compile-clean; visually proofed.
% Uncertainties: the scale-factor, quasitranslation-vector, rotation-matrix,
%                signature, determinant, and continuity-equation
%                modernizations are documented below.
%
% MODERNIZATION CHECK: R_source(t) -> a(t) everywhere. The source's arbitrary
% quasitranslation vector a_source -> bold b in (14.2.7), avoiding a collision
% with the scale factor. Its rotation matrix R -> O, reserving R for curvature.
% Equation (14.2.1) is written as ds^2 in (-+++), and (14.2.16) as -g>0.
% Equation (14.2.19) displays both Weinberg's form and the equivalent modern
% continuity equation, defining H=dot a/a for later use.
% LITERAL-OPERATOR AUDIT: the plus signs in (14.2.7), (14.2.12), and every
% continuation line of (14.2.18), together with the minus sign in 1-kr^2,
% were checked directly against physical PDF pp. 434--436.

\subsection{The Robertson--Walker Metric}\label{sec:14.2}

The formulation of the Cosmological Principle given in the last section
allows us to apply the results of Section~\ref{sec:13.5} for spaces with
maximally symmetric subspaces. We see immediately that it must be possible to
choose coordinates \(t,r,\theta,\phi\), for which the metric takes the form
given in Eq.~\eqref{eq:13.5.32}:
\begin{equation}
  \dd s^2
  =
  -\dd t^2
  a^2(t)
  \left[
    \frac{\dd r^2}{1-kr^2}
    +r^2\dd\theta^2
    +r^2\sin^2\theta\,\dd\phi^2
  \right],
  \tag{14.2.1}\label{eq:14.2.1}
\end{equation}
where \(a(t)\) is an unknown function of time, and \(k\) is a constant, which
by a suitable choice of units for \(r\) can be chosen to have the value
\(+1\), \(0\), or \(-1\). (These are not necessarily the same as the cosmic
standard coordinates introduced in the last section, although \(t\) in
Eq.~\eqref{eq:14.2.1} is the cosmic standard time, or a function of it.) The
metric \eqref{eq:14.2.1} is known in cosmology as the Robertson--Walker
metric.

It is interesting to consider the geometrical properties of the
three-dimensional spaces of constant \(t\). These have metric
\begin{equation}
  {}^{(3)}g_{rr}=\frac{a^2(t)}{1-kr^2},
  \qquad
  {}^{(3)}g_{\theta\theta}=r^2a^2(t),
  \qquad
  {}^{(3)}g_{\phi\phi}=r^2\sin^2\theta\,a^2(t),
  \tag{14.2.2}\label{eq:14.2.2}
\end{equation}
with \({}^{(3)}g_{\mu\nu}\) vanishing for \(\mu\neq\nu\). Comparing with
Eqs.~\eqref{eq:13.3.23}--\eqref{eq:13.3.25} shows that the
three-dimensional curvature constant is
\begin{equation}
  {}^{(3)}K(t)=\frac{k}{a^2(t)}.
  \tag{14.2.3}\label{eq:14.2.3}
\end{equation}
For \(k=-1\) or \(k=0\) the space is infinite, while for \(k=+1\) it is
finite (though unbounded), with proper circumference given by
Eq.~\eqref{eq:13.3.33} as
\begin{equation}
  {}^{(3)}L=2\pi a(t)
  \tag{14.2.4}\label{eq:14.2.4}
\end{equation}
and proper volume given by Eq.~\eqref{eq:13.3.29} as
\begin{equation}
  {}^{(3)}V=2\pi^2a^3(t).
  \tag{14.2.5}\label{eq:14.2.5}
\end{equation}
For \(k=+1\) the spatial universe can be regarded as the surface of a sphere
of radius \(a(t)\) in four-dimensional Euclidean space (see
Section~\ref{sec:13.3}), and \(a(t)\) can justly be called the ``radius of
the universe.'' For \(k=-1\) and \(k=0\) no such interpretation is possible,
but \(a(t)\) still sets the scale of the geometry of space, so \(a(t)\) will
in all cases be called the \emph{cosmic scale factor}.

The construction in Section~\ref{sec:13.5} of the coordinates
\(r,\theta,\phi,t\) was carried out in such a way that the coordinate
transformations which leave the four-dimensional metric
\eqref{eq:14.2.1} form-invariant are just the purely spatial transformations
that leave \eqref{eq:14.2.2} form-invariant. These include the rigid rotations
\begin{equation}
  x'^i=O^i{}_j x^j,
  \qquad i,j=1,2,3,
  \tag{14.2.6}\label{eq:14.2.6}
\end{equation}
where \(O\) is an arbitrary orthogonal matrix. As usual,
\[
  x^1=r\sin\theta\cos\phi,\qquad
  x^2=r\sin\theta\sin\phi,\qquad
  x^3=r\cos\theta.
\]
There are also ``quasitranslations,'' obtained by setting \(KC\) equal to
\(k\) times the unit matrix in Eq.~\eqref{eq:13.3.17}:
\begin{equation}
  \mathbf{x}'
  =
  \mathbf{x}
  +\mathbf{b}
  \left[
    \sqrt{1-k\mathbf{x}^2}
    -
    \left(1-\sqrt{1-k\mathbf{b}^2}\right)
    \frac{\mathbf{x}\cdot\mathbf{b}}{\mathbf{b}^2}
  \right],
  \tag{14.2.7}\label{eq:14.2.7}
\end{equation}
where \(\mathbf{b}\) is an arbitrary three-vector.

The transformation \eqref{eq:14.2.7} carries the origin into the point
\(\mathbf{b}\), so we can conclude that any fixed point can serve as the
origin of a system of coordinates equivalent to the coordinate system of
Eq.~\eqref{eq:14.2.1}. That is, the ``fundamental trajectories'' of
observers, to whom the universe looks the same as it does to us, are just
\(\mathbf{X}(t;\mathbf{b})=\mathbf{b}\). We have already noted in the last
section that the fundamental trajectories ought to be close to the paths of
``typical'' galaxies, so we can tentatively conclude that the spatial
coordinates \(r,\theta,\phi\) form a \emph{comoving system}, in the sense
that typical galaxies have constant spatial coordinates \(r,\theta,\phi\).
One can imagine the comoving coordinate mesh to be like lines painted on the
surface of a balloon, on which dots represent typical galaxies. As the
balloon is inflated or deflated the dots will move, but the lines will move
with them, so each dot will keep the same coordinates.

It is important to note that the fundamental trajectories
\(\mathbf{x}=\text{const}\) are geodesics, because Eq.~\eqref{eq:14.2.1}
gives
\begin{equation}
  \Gamma^\mu{}_{00}=0.
  \tag{14.2.8}\label{eq:14.2.8}
\end{equation}
Thus the statement that a galaxy has constant \(r,\theta,\phi\) is perfectly
consistent with the supposition that galaxies are in free fall. Note also
that the time coordinate \(t\) in \eqref{eq:14.2.1} is not only a possible
``cosmic standard'' time in the sense of the last section; it is also the
proper time told by a clock at rest in any typical freely falling galaxy.
The coordinates \(\mathbf{x},t\) are thus comoving in precisely the same
sense as the ``Gaussian normal'' coordinates introduced in
Section~\ref{sec:11.8}.

We can obtain a deeper insight into the behavior of matter in a
Robertson--Walker universe by applying the Cosmological Principle to the
tensors that describe the average state of cosmic matter, such as the
energy-momentum tensor \(T^{\mu\nu}\), and the current \(J_G^\mu\) of
galaxies. (\(J_G^\mu\) is defined exactly like the electric current
\eqref{eq:5.2.13}, but the sum runs over galaxies instead of particles, and
a factor \(1\) replaces \(e_n\).) All such tensors are required to be
form-invariant [in the sense of Section~\ref{sec:13.4} or
Eq.~\eqref{eq:14.1.2}] with respect to those coordinate transformations,
such as \eqref{eq:14.2.6} and \eqref{eq:14.2.7}, which leave the metric
\eqref{eq:14.2.1} form-invariant. These ``isometries'' are purely spatial,
so they transform \(J_G^0\) and \(T^{00}\) as three-scalars;
\(J_G^i\) and \(T^{i0}\) as three-vectors; and \(T^{ij}\) as a
three-tensor. According to the theorems proved in Section~\ref{sec:13.4},
this then requires that
\begin{equation}
  J_G^0=n_G(t),
  \qquad
  J_G^i=0,
  \tag{14.2.9}\label{eq:14.2.9}
\end{equation}
and
\begin{equation}
  T_{00}=\rho(t),
  \qquad
  T_{0\ii}=0,
  \qquad
  T_{ij}={}^{(3)}g_{ij}p(t),
  \tag{14.2.10}\label{eq:14.2.10}
\end{equation}
where \(n_G,\rho,\) and \(p\) are unknown quantities that may depend on
\(t\), but not on \(r,\theta,\) or \(\phi\). These results can be written
more elegantly as
\begin{equation}
  J_G^\mu=n_GU^\mu,
  \tag{14.2.11}\label{eq:14.2.11}
\end{equation}
\begin{equation}
  T_{\mu\nu}
  =
  (\rho+p)U_\mu U_\nu
  p\,g_{\mu\nu},
  \tag{14.2.12}\label{eq:14.2.12}
\end{equation}
where \(U^\mu\) is a ``velocity four-vector''
\begin{equation}
  U^0=1,
  \tag{14.2.13}\label{eq:14.2.13}
\end{equation}
\begin{equation}
  U^i=0.
  \tag{14.2.14}\label{eq:14.2.14}
\end{equation}
Equation~\eqref{eq:14.2.14} shows that \emph{the contents of the universe
are, on the average, at rest in the coordinate system \(r,\theta,\phi\)}, as
expected. In addition, comparison of Eq.~\eqref{eq:14.2.12} with
Eq.~\eqref{eq:5.4.2} shows that \emph{the energy-momentum tensor of the
universe necessarily takes the same form as for a perfect fluid}.

It will be useful to have on hand the differential equations for
\(n_G(t)\), \(\rho(t)\), and \(p(t)\) that are provided by conservation
principles. If galaxies are neither created nor destroyed, then \(J_G^\mu\)
obeys the conservation equation \eqref{eq:5.2.14}:
\begin{equation}
  0
  =
  \nabla_\mu J_G^\mu
  =
  \frac{1}{\sqrt{-g}}
  \partial_\mu\!\left(\sqrt{-g}\,J_G^\mu\right)
  =
  \frac{1}{\sqrt{-g}}
  \frac{\dd}{\dd t}\!\left(\sqrt{-g}\,n_G\right).
  \tag{14.2.15}\label{eq:14.2.15}
\end{equation}
The metric \eqref{eq:14.2.1} has \(g=\det(g_{\mu\nu})\) given by
\begin{equation}
  -g
  =
  \frac{a^6(t)r^4\sin^2\theta}{1-kr^2},
  \tag{14.2.16}\label{eq:14.2.16}
\end{equation}
and therefore the conservation of galaxies yields the relation
\begin{equation}
  n_G(t)a^3(t)=\text{constant}.
  \tag{14.2.17}\label{eq:14.2.17}
\end{equation}
(Note that \(n_G\) is the number density per unit proper volume, and
therefore increases or decreases according as the universe shrinks or
expands, while \(n_Ga^3\) is the number density per unit coordinate volume,
and therefore remains constant in a comoving coordinate system.) The
energy-momentum tensor \eqref{eq:14.2.12} obeys the conservation equation
\eqref{eq:5.4.3}:
\begin{align}
  0
  ={}&
  \nabla_\nu T^{\mu\nu}
  \notag\\
  ={}&
  \partial^\mu p
  +\frac{1}{\sqrt{-g}}
   \partial_\nu\!\left[
     \sqrt{-g}\,(\rho+p)U^\mu U^\nu
   \right]
  +\Gamma^\mu{}_{\nu\lambda}
   (\rho+p)U^\nu U^\lambda.
  \tag{14.2.18}\label{eq:14.2.18}
\end{align}
Using Eqs.~\eqref{eq:14.2.8} and \eqref{eq:14.2.14}, we find that this
equation is trivially satisfied for \(\mu=r,\theta,\phi\), while for
\(\mu=0\) it reads
\begin{equation}
  a^3(t)\frac{\dd p(t)}{\dd t}
  =
  \frac{\dd}{\dd t}
  \left\{
    a^3(t)\bigl[\rho(t)+p(t)\bigr]
  \right\}
  \quad\Longleftrightarrow\quad
  \dot\rho+3H(\rho+p)=0,
  \qquad
  H\equiv\frac{\dot a}{a}.
  \tag{14.2.19}\label{eq:14.2.19}
\end{equation}
For instance, if the pressure of cosmic matter is negligible, then
\eqref{eq:14.2.19} gives a result analogous to
Eq.~\eqref{eq:14.2.17}:
\begin{equation}
  \rho(t)a^3(t)=\text{constant}.
  \tag{14.2.20}\label{eq:14.2.20}
\end{equation}

The very great convenience of a comoving coordinate system should not blind
us to the fact that the typical galaxies actually do move further apart or
closer together when \(a(t)\) increases or decreases. To see this clearly,
we need to consider what we mean by the distance between galaxies. Imagine a
chain of typical galaxies lying close together on the line of sight between
us and a distant galaxy at \(r_1,\theta_1,\phi_1\), and suppose that
\emph{at the same cosmic time \(t\)}, observers in each galaxy measured the
distance to the next galaxy, say by measuring the travel time for light
signals. (Note that this is not the same as measuring the time for a single
light signal to go from \(r=0\) to \(r=r_1\).) Adding up all these
subdistances gives the \emph{proper distance}
\begin{equation}
  d_{\mathrm{prop}}(t)
  =
  \int_0^{r_1}\sqrt{g_{rr}}\,\dd r
  =
  a(t)\int_0^{r_1}\frac{\dd r}{\sqrt{1-kr^2}}.
  \tag{14.2.21}\label{eq:14.2.21}
\end{equation}
\footnote{Obviously, no one is going to organize this sort of cosmic
conspiracy, so the proper distance is not very relevant to observational
cosmology. However, we shall see in Section~\ref{sec:14.4} that the more
relevant measures of distance, based on apparent luminosities and angular
diameters, all approach the proper distance \eqref{eq:14.2.21} for
\(r_1\ll1\). Thus, in one sense or another, galaxies do move apart when
\(a(t)\) increases, or together when \(a(t)\) decreases.}

Cosmological theory presents observational astronomy with the challenge of
measuring the function \(a(t)\) and determining whether \(k\) is \(+1\),
\(0\), or \(-1\). This is not all there is to cosmology, but it is a central
problem that must be solved if we are to understand the universe. The balance
of this chapter will describe how well this challenge has been met.
